\section{Záver}

Cieľom tejto práce bolo analyzovať princípy Clean Code, ich praktické využitie a vzťah k moderným nástrojom umelej inteligencie používaným pri vývoji softvéru. Ukázalo sa, že Clean Code predstavuje súbor odporúčaní zameraných na zlepšenie čitateľnosti, testovateľnosti a udržiavateľnosti kódu, avšak jeho aplikácia by nemala byť dogmatická. Nadmerná abstrakcia, predčasná optimalizácia alebo zbytočná architektonická komplexita môžu viesť k vyšším nákladom na vývoj a údržbu bez primeraného prínosu.

Práca zároveň poukázala na to, že nástroje umelej inteligencie dokážu významne urýchliť tvorbu zdrojového kódu, generovanie testov a refaktoring existujúcich riešení. Kvalita výsledného kódu však nezávisí iba od schopností modelu, ale predovšetkým od kvality vstupného kontextu, spôsobu integrácie do vývojového procesu a úrovne ľudského dohľadu. AI je schopná generovať kód, ktorý spĺňa mnohé princípy Clean Code, no zároveň môže vytvárať riešenia obsahujúce architektonické nedostatky, technický dlh alebo logické chyby.

Porovnanie čistého a nekonzistentného AI generovaného kódu ukazuje, že krátkodobé zrýchlenie vývoja nemusí automaticky znamenať dlhodobú úsporu nákladov. Zatiaľ čo kvalitne kontrolovaný AI generovaný kód dosahuje podobné náklady na údržbu ako tradičný vývoj, nekvalitný kód môže v horizonte niekoľkých rokov výrazne zvýšiť technický dlh a náklady na údržbu systému.

Na základe získaných poznatkov možno konštatovať, že najefektívnejším prístupom je kombinácia princípov Clean Code, automatizovaných kontrol kvality, testovania a pravidelného code review. Umelá inteligencia by mala byť vnímaná ako podporný nástroj, ktorý zvyšuje produktivitu vývojárov, nie ako náhrada za odborné rozhodovanie. Dlhodobú kvalitu softvéru zabezpečuje predovšetkým schopnosť tímu rozumieť vytvorenému kódu, priebežne merať jeho kvalitu a robiť pragmatické rozhodnutia založené na reálnych potrebách projektu.